\documentclass[10pt,a4paper]{amsart}
\usepackage[margin=19mm]{geometry}
\usepackage{amsmath,amssymb,mathtools}
\usepackage[scr=rsfs,cal=euler]{mathalfa}
\usepackage{xfrac}
\usepackage[unicode,hidelinks,bookmarks=false]{hyperref}
\newcommand{\ie}{i.e.\ }
\newcommand{\cf}{cf.\ }
\newcommand{\resp}{respectively\ }
\newcommand{\Subsection}{\subsection}
\providecommand{\nobreakdash}{\nobreak\hbox{-}}
\setlength{\emergencystretch}{2em}
\multlinegap0pt
\usepackage{bm}
\usepackage{bbm}
\usepackage{dsfont}
\usepackage{xspace}
\usepackage{cancel}
\usepackage{mathtools}
\usepackage{enumitem}
\usepackage{amsthm}
\makeatletter
\@ifundefined{theorem}{\newtheorem{theorem}{Theorem}}{}
\@ifundefined{lemma}{\newtheorem{lemma}[theorem]{Lemma}}{}
\makeatother
\newcommand{\Ins}[1]{{\rm Int}{(#1\cup \substrate\cup \Sigma)}}
\newcommand{\Lip}{\mathrm{Lip}}
\newcommand{\R}{\mathbb{R}}
\newcommand{\bound}{b}
\newcommand{\cH}{\mathcal{H}}
\newcommand{\cl}[1]{\overline{#1}}
\newcommand{\p}{\partial}
\newcommand{\strictlyincluded}{\subset\subset}
\newcommand{\substrate}{S}
\newcommand{\supp}{\mathrm{supp}\,}
\renewcommand{\S}{\mathbb{S}}
\renewcommand{\d}{\mathrm{d}}
\title[Existence of minimizers for the SDRI model in $\mathbb{R}^n$]{Existence of minimizers for the SDRI model in $\mathbb{R}^n$: Wetting and dewetting regimes with mismatch strain}
\author{Shokhrukh Kholmatov, Paolo Piovano}
\date{Source: arXiv:2305.10304v2 (CC BY 4.0)}
\begin{document}
\maketitle

\noindent\emph{Source and licence.} Existence of minimizers for the SDRI model in $\mathbb{R}^n$: Wetting and dewetting regimes with mismatch strain. Version arXiv:2305.10304v2 (CC BY 4.0), distributed under the Creative Commons Attribution 4.0 International licence, as recorded in the arXiv licence metadata for this exact version and shown on the abstract page https://arxiv.org/abs/2305.10304v2. Full rights notice: licenses/arxiv-2305.10304-CC-BY-4.0.txt.\par\smallskip

\noindent\textit{Editorial scope.} Counted unit: the Lemma whose source label is \texttt{lem:holes\_near\_Sigma} in the article (source0.tex lines 1149--1167, source label \texttt{lem:holes\_near\_Sigma}), delivered together with the article's own complete original proof of it (source0.tex lines 1169--1293, one \texttt{proof} environment of 125 lines). No mathematics is written, repaired or re-derived: statement and proof are the article's own text line for line. Required context retained uncounted: fjskiii098 (lines 872--874). Every definition, display and label of those lines is retained unchanged. Omitted from this package and not redistributed: 1--871, 875--1148, 1168--1168, 1294--4229. Nothing inside a retained excerpt is omitted or altered: every retained body is delivered exactly as the article's lines are written, so no omission receipt of any kind accompanies this package. Citations: the retained excerpts carry no citation command at all, so no bibliography entry of the article is transcribed and no reference is redistributed. Reference adaptation: none was needed and none was made; every reference inside the delivered unit resolves inside it, so no unresolved reference or citation is produced. Printed slips kept literal: no printed word, symbol, hypothesis or number of the counted statement or of its proof was altered. Layout and class: the article's own class is \texttt{amsart} with packages including amssymb, amsmath, amsthm, amsfonts, color, geometry, enumitem, comment, graphicx, setspace, wrapfig, hyperref; the wrapper is the lane's pinned \texttt{amsart} editorial wrapper at 10pt on A4, which restates the theorem-like environments the retained text needs and adds the article's own macro definitions that the retained text needs (\texttt{\textbackslash Ins}, \texttt{\textbackslash Lip}, \texttt{\textbackslash R}, \texttt{\textbackslash S}, \texttt{\textbackslash bound}, \texttt{\textbackslash cH}, \texttt{\textbackslash cl}, \texttt{\textbackslash d}, \texttt{\textbackslash p}, \texttt{\textbackslash strictlyincluded}, \texttt{\textbackslash substrate}, \texttt{\textbackslash supp}), with extra wrapper packages (bm, bbm, dsfont, xspace, cancel, mathtools, enumitem). The delivered numbering is the unit's own. Assets: no retained span contains a figure environment, a graphics-inclusion command, a PDF-page inclusion, a table or a TikZ picture; no image file is copied into this package, the package declares no asset dependency and no image file is redistributed with it. Licence: version arXiv:2305.10304v2 of this article is distributed under the Creative Commons Attribution 4.0 International licence, as recorded in the arXiv licence metadata for this exact version and shown on the abstract page https://arxiv.org/abs/2305.10304v2. A full rights notice is delivered at licenses/arxiv-2305.10304-CC-BY-4.0.txt. Characters: every retained excerpt is pure ASCII, so the delivered engine, XeLaTeX with its default Latin Modern text font, reports no missing character.
 \begin{equation}\label{fjskiii098}
\bound_1\le \phi(\nu) \le \bound_2,\quad \nu\in\S^{n-1},
\end{equation}

\begin{lemma}\label{lem:holes_near_Sigma}
Let $U\strictlyincluded \Ins{\Omega}$ be an open set, $\delta\in(0,1)$ and  $K\subset U\cap\Sigma$ be any $\cH^{n-1}$-measurable set. Then there exist an open set $C\subset U\cap \Omega$ of finite perimeter such that 
\begin{itemize}[left=15pt]
 \item[\rm(i)] $C\strictlyincluded U$ and $\cH^{n-1}(\p C\setminus \p^*C)=0;$
 
 \item[\rm(ii)] $\cH^{n-1}(K\setminus \p C)=\cH^{n-1}(K\setminus \cl{C})<\delta$ and $|C|<\delta;$
 
 \item[\rm(iii)] $\cH^{n-1}(U\cap \Sigma \cap\p C \setminus K)<\delta;$
 
 \item[\rm(iv)] for any norm $\phi$ in $\R^n$ satisfying \eqref{fjskiii098}
 $$
 \int_{\Omega\cap \p C} \phi(\nu_C)\d\cH^{n-1} \le \int_K\phi(\nu_\Sigma) \d\cH^{n-1} +\delta, 
 $$
 and 
 $$
 \int_{\p C} \phi(\nu_C)\d\cH^{n-1} \le 2\int_K\phi(\nu_\Sigma) \d\cH^{n-1} +\delta .
 $$ 
\end{itemize}
\end{lemma}

\begin{proof}
Let 
$$
\epsilon:= \frac{\delta}{8(1+\bound_2)(1 + \cH^{n-1}(\Sigma))}.
$$
We divide the proof into two steps.

{\it Step 1.} Let $Q_r(x_0)\subset U$ be a cube centered at $x\in\Sigma$ such that $\Sigma\cap Q_r(x_0) = {\rm graph}(f)$ for some Lipschitz function $f:V\to\R$ and a cube $V\subset \R^{n-1},$ and assume that $S\cap Q_r(x_0)$ is a subgraph of $f.$ 
Let $V''\strictlyincluded V'\strictlyincluded V$ be open sets such that 
$$
\cH^{n-1}(V\setminus V'') < \frac{\cH^{n-1}(Q_r(x_0)\cap\Sigma)}{1 + \Lip(f)}\,\epsilon
$$ 
and for $\gamma\in(0,\frac{\cH^{n-1}(Q_r(x_0)\cap\Sigma)}{1+\cH^{n-1}(V)}\,\epsilon)$ let $g\in \Lip_c(V;[0,\gamma])$ be such that $g\equiv \gamma$ in $V'',$ $\supp(g)=V'$ and $\Lip(g)<1.$ We may assume that $\gamma$ is so small that the set open set $C,$ whose boundary lies on the graphs of $f$ and $f+g,$ is compactly contained in $Q_r(x_0)$ and $C\cap S=\emptyset.$
Then 
\begin{align*}
\int_{\Omega\cap \p C}  \phi(\nu_C)\d\cH^{n-1} = & \int_{V'} \phi(-\nabla(f+g),1)\d\cH^{n-1} \le 
\int_{V'} \Big(\phi(-\nabla f,1) + \phi(-\nabla g,0)\Big)\d\cH^{n-1} \\
\le & \int_{V} \phi(-\nabla f,1)\d\cH^{n-1} + \bound_2\Lip(f)\cH^{n-1}(V'\setminus V'')\\
< & \int_{Q_r(x_0)\cap\Sigma}\phi(\nu_\Sigma)\d\cH^{n-1} + \bound_2\cH^{n-1}(Q_r(x_0)\cap\Sigma)\epsilon.
\end{align*}
Similarly, 
\begin{align*}
\int_{\p C}  \phi(\nu_C)\d\cH^{n-1} = & \int_{V'} [\phi(-\nabla(f+g),1) + \phi(-\nabla f,1)]\d\cH^{n-1} \\
\le &  
\int_{V'} \Big(2\phi(-\nabla f,1) + \phi(-\nabla g,0)\Big)\d\cH^{n-1} \\
< & 2\int_{Q_r(x_0)\cap\Sigma}\phi(\nu_\Sigma)\d\cH^{n-1} + \bound_2\cH^{n-1}(Q_r(x_0)\cap\Sigma)\epsilon.
\end{align*}
%
Also by the Fubini's theorem,
$$
|C| = \int_{V'} g \d x'\le \gamma \cH^{n-1}(V') < \cH^{n-1}(Q_r(x_0)\cap\Sigma)\,\epsilon.
$$
Finally, 
$$
\begin{aligned}
\cH^{n-1}(Q_r(x_0)\cap\Sigma\setminus \p C) = & \cH^{n-1}(Q_r(x_0)\cap\Sigma\setminus \cl{C}) 
=\int_{V\setminus V'}\sqrt{1+|\nabla f|^2}\d \cH^{n-1}\\
\le & (1+\Lip(f)) \cH^{n-1}(V\setminus V') < \cH^{n-1}(Q_r(x_0)\cap\Sigma)\,\epsilon.
\end{aligned}
$$

{\it Step 2.} Since $\Sigma$ is Lipschitz and $K$ is $\cH^{n-1}$-rectifiable, we can find a finite  family $Q_{r_1,\nu_1}(x_1),\ldots, Q_{r_m,\nu_m}(x_m) \subset U$ of pairwise disjoint cubes centered at $K$ such that 
\begin{itemize}
 \item[($\rm a_1$)] for each $j,$ $\Sigma\cap Q_{r_j,\nu_j}(x_j)$ is a graph of a Lipschitz function in $\nu_j$ direction;

 \item[($\rm a_2$)] $\theta(K,x_j)=\theta(\Sigma,x_j)=1,$ and the unit normals $\nu_K(x_j)$ and $\nu_{\Sigma}(x_j)$ exist and coincide with $\nu_j;$
 
 \item[($\rm a_3$)] $(1-\epsilon)r_j^{n-1} <\cH^{n-1}(Q_{r_j,\nu_j}(x_j)\cap \Sigma\cap K)\le \cH^{n-1}(Q_{r_j,\nu_j}(x_j)\cap \Sigma)<(1+\epsilon)r_j^{n-1};$
 
 \item[($\rm a_4$)] $\cH^{n-1}\Big(K\setminus \bigcup_{j=1}^m Q_{r_j,\nu_j}(x_j)\Big)<\epsilon.$
\end{itemize}

Note that by ($\rm a_3$),
\begin{equation}\label{K_out_sigma89}
\cH^{n-1}(Q_{r_j,\nu_j}(x_j)\cap \Sigma\setminus K) <2\epsilon r_j^{n-1} < \frac{2\epsilon}{1-\epsilon} \cH^{n-1}(Q_{r_j,\nu_j}(x_j)\cap \Sigma). 
\end{equation}
%
By Step 1, for each $j$ we can contruct an open set $C_j\strictlyincluded Q_{r_j,\nu_j}(x_j)$ satisfying $C_j\cap\substrate =\emptyset,$
\begin{align}\label{est_Cjjjj9099}
\int_{\Omega\cap \p C_j} \phi(\nu_{C_j})\d\cH^{n-1} < \int_{Q_{r_j,\nu_j}(x_j)\cap\Sigma}\phi(\nu_\Sigma)\d\cH^{n-1} + \bound_2\cH^{n-1}(Q_{r_j,\nu_j}(x_j)\cap\Sigma)\epsilon 
\end{align}
and 
\begin{align*} 
\int_{\p C_j} \phi(\nu_{C_j})\d\cH^{n-1} < 2 \int_{Q_{r_j,\nu_j}(x_j)\cap\Sigma}\phi(\nu_\Sigma)\d\cH^{n-1} + \bound_2\cH^{n-1}(Q_{r_j,\nu_j}(x_j)\cap\Sigma)\epsilon.
\end{align*}
Moreover,
\begin{equation}\label{vol_estCjjjj}
|C_j| < \cH^{n-1}(Q_{r_j,\nu_j}(x_j)\cap\Sigma)\epsilon  
\end{equation}
%
and
%
\begin{align}\label{out_jshsh00}
\cH^{n-1}(Q_{r_j,\nu_j}(x_j)\cap\Sigma\setminus \p C_j) < \cH^{n-1}(Q_{r_j,\nu_j}(x_j)\cap\Sigma)\,\epsilon.
\end{align}


We claim that $C= \bigcup_{i=1}^m C_j$ satisfies the assertions of the lemma.

(i) By construction, $C\strictlyincluded U$ and since each $C_j$ is almost Lipschitz, $\cH^{n-1}(\p C_j\setminus \p^*C_j)=0.$ Hence, by the pairwise disjointness of $\cl{C_j},$ $\cH^{n-1}(\p C\setminus \p^*C)=0$ and (i) follows.

(ii) By ($\rm a_4$) and \eqref{out_jshsh00},
\begin{align*}
\cH^{n-1}(K\setminus \p C) \le & \cH^{n-1}\Big(K\setminus \bigcup_{j=1}^m Q_{r_j,\nu_j}(x_j)\Big) + 
\sum\limits_{j=1}^m \cH^{n-1}( Q_{r_j,\nu_j}(x_j)\cap \Sigma  \setminus \p C_j) \\
< & \epsilon +  \sum\limits_{j=1}^m \cH^{n-1}( Q_{r_j,\nu_j}(x_j) \cap\Sigma)\,\epsilon \le
(1+\cH^{n-1}(\Sigma))\epsilon <\delta.
\end{align*}
Moreover, by \eqref{vol_estCjjjj},
$$
|C| \le \sum\limits_{j=1}^n |C_j| \le \cH^{n-1}(\Sigma)\epsilon <\delta.
$$


(iii) By \eqref{K_out_sigma89},
\begin{align*}
\cH^{n-1}( U\cap\Sigma\cap \p C \setminus K) = & \sum\limits_{j=1}^m \cH^{n-1}(Q_{r_j,\nu_j}(x_j)\cap \Sigma\cap \p C_j \setminus K) \\
\le & \sum\limits_{j=1}^m \cH^{n-1}(Q_{r_j,\nu_j}(x_j)\cap \Sigma\setminus K) <\frac{2\epsilon}{1-\epsilon}\cH^{n-1}(\Sigma) <\delta.
\end{align*}

(iv) Since $\p^* C\subset \cup_j \p^* C_j,$ by \eqref{est_Cjjjj9099} we have
$$
\int_{\Omega\cap\p C} \phi(\nu_C)\d\cH^{n-1} 
\le \sum\limits_{j=1}^m \int_{Q_{r_j,\nu_j}(x_j)\cap\Sigma}\phi(\nu_\Sigma)\d\cH^{n-1} + \bound_2\cH^{n-1}(\Sigma)\epsilon
$$
Moreover, by \eqref{fjskiii098},
$$
\int_{Q_{r_j,\nu_j}(x_j)\cap\Sigma}\phi(\nu_\Sigma)\d\cH^{n-1} \le 
\int_{Q_{r_j,\nu_j}(x_j)\cap K} \phi(\nu_\Sigma)\d\cH^{n-1} + \bound_2 \cH^{n-1}(Q_{r_j,\nu_j}(x_j)\cap \Sigma\setminus K),
$$
and thus by \eqref{K_out_sigma89},
$$
\int_{\Omega\cap\p C} \phi(\nu_C)\d\cH^{n-1}  
\le 
\int_K \phi(\nu_\Sigma)\d\cH^{n-1} + \Big(\tfrac{\bound_2}{1-\epsilon} + \bound_2) \cH^{n-1}(\Sigma)\epsilon < \int_K \phi(\nu_\Sigma)\d\cH^{n-1} +\delta.
$$
Finally, since $\Sigma\cap \p C\subset K\cup (\Sigma\cap \p C\setminus K),$  
\begin{align*}
\int_{\p C} \phi(\nu_C)\d\cH^{n-1} = & \int_{\Omega\cap\p C} \phi(\nu_C)\d\cH^{n-1} + \int_{\Sigma\cap \p C} \phi(\nu_\Sigma)\d\cH^{n-1}\\
\le & 2\int_K\phi(\nu_\Sigma)\d\cH^{n-1} + 3\bound_2\cH^{n-1}(\Sigma) \epsilon + 
\bound_2\cH^{n-1}(\Sigma\cap \p C\setminus K)\\
< & 2\int_K\phi(\nu_\Sigma)\d\cH^{n-1} + 7\bound_2\cH^{n-1}(\Sigma) \epsilon 
< 2\int_K \phi(\nu_\Sigma)\d\cH^{n-1} +\delta. 
\end{align*}
\end{proof}

\end{document}
